% Generated by tools/edn_latex.py from reports/edn.md.
% Do not edit: change reports/edn.md and run `make edn`.
\ednentry{
  id = {EDN-06},
  title = {Success criteria for the price model},
  date = {2026-09-22},
  milestone = {M1: Project Inception},
  topic = {Performance Criteria},
  participants = {@lukas2510},
  decision = {A model is deployable when it meets SC-01 to SC-05 on the grouped test set: MdAPE \ensuremath{\leq} 9\,\%; at least 85\,\% of predictions within \ensuremath{\pm}20\,\%; MdAPE at least 30\,\% lower than the median baseline B0; MdAPE \ensuremath{\leq} 15\,\% in every input-based segment with at least 500 test rows; empirical coverage of nominal 90\,\% intervals between 88\,\% and 92\,\% for full inputs and for a partial-input scenario.},
  rationale = {\emph{Alternatives considered:}
\begin{itemize}[leftmargin=*, topsep=2pt]
\item \textbf{Option A (chosen): absolute targets plus a baseline comparison, per segment.}
\newline Pros: says whether the product is useful to users and whether the model earns its complexity; segment criterion catches models that are good on average but bad for a group; grounded in an exploratory run.
\newline Cons: targets depend on this dataset and must be revisited if the scope changes.
\item \textbf{Option B: baseline comparison only.}
\newline Pros: independent of the data.
\newline Cons: says nothing about usefulness; a weak baseline makes it easy to pass.
\item \textbf{Option C: MAPE-based targets, as in the original plan.}
\newline Pros: common and familiar.
\newline Cons: dominated by cheap cars and outliers (the same model has 9.7\,\% MAPE but 6.8\,\% MdAPE).
\end{itemize}
\par
\emph{Why:} An exploratory run on the final scope (102,695 listings, 80/20 split grouped by seller) gave B0 11.6\,\% MdAPE / 71.4\,\% within \ensuremath{\pm}20\,\% and an untuned LightGBM 6.8\,\% / 90.7\,\%. The targets keep a margin to the LightGBM values so they hold across splits, while ruling out a model only marginally better than B0. Price buckets are excluded from the segment criterion because they condition on the target. The run already shows one gap: cars older than 20 years reach 16.7\,\% MdAPE, which is recorded as a known risk rather than hidden by loosening SC-04. Re-checked after the \texttt{ES} holdout (EDN-03) on 96,831 listings: B0 11.9\,\% / 70.9\,\%, LightGBM 6.7\,\% / 90.6\,\%, cars older than 20 years 15.3\,\%; the criteria stay unchanged.},
  ai = {Information seeking, Alternative generation, Alternative assessment, Recommendation},
  response = {Accepted},
  assessment = {AI ran the exploratory baseline and LightGBM models, proposed the criteria with concrete thresholds and recommended option A. Lukas accepted it because the thresholds are tied to measured values. After the decision, AI's per-segment check found that the proposed SC-04 is not yet met for cars older than 20 years; the criterion was kept as agreed and the gap documented.},
  aievidence = {Claude Code session on 2026-09-22 while working on issue \#2: prompt ``Which kind of success criteria?'' with options A to C and the exploratory results; Lukas chose option A.},
  evidence = {\href{https://github.com/mlops-2627q1-mds-upc/MLOPS_Recommenditos/blob/main/docs/docs/problem-spec.md}{Problem specification}, sections 6 to 8; \href{https://github.com/mlops-2627q1-mds-upc/MLOPS_Recommenditos/issues/2}{issue \#2}, \href{https://github.com/mlops-2627q1-mds-upc/MLOPS_Recommenditos/pull/17}{PR \#17}.},
}
